\documentclass{article}
\title{チャットサーバー}
\author{
23ts028 河村 貴史\\
23ts030 川瀬 翔大
}
\date{}
\begin{document}
\maketitle
\section{全体の構造}
本システムは、C言語で実装されたクライアント/サーバー型のチャットアプリケーションです。基本的な通信プロトコルとしてTCPソケット通信を採用し、サーバーとクライアント間でリアルタイムなデータ交換を実現しています。
\subsection{サーバー(server.c)}
サーバーは複数のクライアント接続を同時に処理するためにマルチスレッドモデルを採用しています。
\begin{itemize}
    \item 同時接続：最大10クライアントの同時接続をサポートし、接続ごとに専用の処理スレッド（handle\_client）を立ち上げます。
    \item ユーザー管理：接続中のクライアントリストはミューテックス（clients\_mutex）によってスレッドセーフに管理されます。
    \item ユーザー情報の永続化：ユーザー認証情報（名前、パスワード）はjanssonライブラリを利用し、JSON形式（users.json）でファイルに保存・ロードされます。
\end{itemize}
\subsection{クライアント(client.c)}
クライアントは、ユーザーからの入力を受け付け、サーバーからのメッセージを受信する機能を担います。
\begin{itemize}
    \item 非同期受信：receive\_messagesという独立したスレッドでサーバーからのメッセージを待ち受けることで、ユーザーが標準入力からメッセージを入力中でも、メッセージの受信がブロックされない設計になっています。
\end{itemize}
\section{登録(ユーザーアカウント作成)}
ユーザー登録はクライアントとサーバー間の以下の手順で行われます。\\\\
\textbf{クライアント操作の流れ}
\begin{enumerate}
    \item モード選択：「1:登録 2:ログイン」のプロンプトに対し、1（登録）を入力し、サーバーに送信します。
    \item 名前入力：「名前を入力」のプロンプトに従い、ユーザー名を送信します。
    \item パスワード入力：「パスワードを入力」のプロンプトに従い、パスワードを送信します。
    \item 結果受信：サーバーからの「登録完了!」メッセージを受信し、認証済み状態としてチャットセッションを開始します。
\end{enumerate}
\textbf{サーバー側の主な処理}
\begin{enumerate}
    \item データ保存：サーバーは、add\_user関数を実行し、新しいユーザー情報をJSON形式でusers.jsonファイルに追記・永続化します。
    \item メモリ更新：load\_users関数によりメモリ上のユーザーリストを最新化します。
    \item 応答：クライアントに対して「登録完了!」メッセージを送信し、認証済み状態としてチャットセッションを開始します。
\end{enumerate}
\section{ログイン}
既存ユーザーのログイン処理は以下の手順で行われます。
\textbf{クライアント操作の流れ}
\begin{enumerate}
    \item モード選択：「1:登録 2:ログイン」のプロンプトに対し、2（ログイン）を入力し、サーバーに送信します。
    \item 名前入力：「名前を入力」のプロンプトに従い、ユーザー名を送信します。
    \item パスワード入力：「パスワードを入力」のプロンプトに従い、パスワードを送信します。
    \item 再試行：認証失敗の場合、「認証失敗、再入力」メッセージを受信し、名前とパスワードの再入力を繰り返します。（クライアント側でのループ制御は課題です）
    \item 接続：認証成功後、「ログイン成功!」メッセージを受信し、チャットに進みます。
\end{enumerate}
\textbf{サーバー側の主な処理}
\begin{enumerate}
    \item 認証情報の送信：クライアントは「名前」と「パスワード」をサーバーに送信します。
    \item 認証：サーバーはload\_usersで最新のユーザーリストを読み込んだ後、authenticate関数を使用して提供された認証情報を照合します。
    \item 再試行制御：認証に失敗した場合、サーバーはクライアントに「認証失敗、再入力」メッセージを送信し、再入力を促します。
    \item 接続：認証が成功すると、クライアントは接続リストに追加され、チャット機能へ進みます。
\end{enumerate}
\section{チャット}
本システムは、メッセージのやり取りに関して「全体トーク」と「個別トーク」の2つのモードを提供します。
\subsection{全体トーク}
\begin{itemize}
    \item クライアントが送信したメッセージは、サーバーで送信者名が付加されて整形されます（例: [ユーザー名] メッセージ）。
    \item broadcast\_message関数が、送信者自身を除く全ての接続中クライアントに対してメッセージを一斉送信します。
\end{itemize}
\begin{itemize}
    \item 個別トークは、メッセージの先頭に特定のコマンド（/whisper）を使用することで実現されます。
    \item コマンド形式: /whisper <宛先ユーザー名> <メッセージ本文>
    \item サーバーはコマンドを解析し、send\_private\_message関数を用いて、指定されたユーザーにのみメッセージを送信します。
\end{itemize}
\section{デプロイ}
Dockerを用いてコンテナ化しており、環境構築の手間を削減して、どの環境でも一貫した動作を保証しています。
\subsection{Dockerfileの構造}
\begin{enumerate}
    \item ベースイメージ：ubuntu22:04を使用し、安定したLinux環境を提供します。
    \item 依存関係のインストール：build-essential（GCC含む）とlibjansson-devをインストールします。
    \item コンパイル：server.cとusers.jsonをコンテナにコピー後、gccコマンドによりサーバーの実行可能ファイルを生成します。
    \item ポート番号：EXPOSE 8080により、外部からアクセス可能なポートとして8080を指定しています。
    \item CMD ["./server"]により、コンテナの起動と同時にサーバーアプリケーションが実行されます。
\end{enumerate}
\section{課題点}
今回チャットアプリケーションを作成したが、まだまだ課題点、修正点が存在する。課題点を以下にまとめる。
\begin{itemize}
    \item 時刻の未表示：メッセージの送受信の時刻の表示をできるようにしたい
    \item トーク履歴の確認：メッセージを一つ１つをjsonファイルに書き込みトーク開始時に今までのトーク履歴の表示を行いたい。
    \item 個別トーク：個別トークを選択したときに、ユーザーの一覧が表示され、そこから、個別トークをできるようにしたかった。現状はコマンドで送信することしかできない。
    \item 通信環境の制約：同じWi-Fi環境でしか通信ができない
    \item コード整形：コードか言葉かを判別してトークは見やすいように整形できるようにしたい。
\end{itemize}
\section{まとめ}
今回チャットアプリケーションの作成を通じて、TCPソケット通信、コンテナの作成、jsonライブラリについて理解が深まった。しかし、様々な課題も見つかった。今後は見つかった課題を一つ１つ解決できるように励みたい。

\begin{thebibliography}{9}
    \bibitem{1}https://daeudaeu.com/socket/,2025年11月25日 \bibitem{2}https://qiita.com/fujifuji1414/items/6daa393a86582d81f0b5,2025年11月25日
    \bibitem{3}https://qiita.com/Sicut\_study/items/4f301d000ecee98e78c9,2025年11月25日
    \bibitem{4}https://zenn.dev/pataro/articles/04dae702e9b073,2025年11月25日
    \bibitem{5}https://af-e.net/c-language-pthread-usage/,2025年11月25日

\end{thebibliography}
\end{document}